\section{Implementation}
\label{sec:implementation}

\todo{Architecture figure: \texttt{figures/architecture.pdf}. Build via
\texttt{make figures}; see \texttt{figures/architecture.mmd}.}

\subsection{Parser and AST}
A tree-sitter grammar (\texttt{src/tree-sitter-yuho/grammar.js}) emits a
concrete syntax tree which a Python AST builder
(\texttt{src/yuho/ast/builder.py}, $\sim$1.4k SLOC) lifts into immutable
frozen-dataclass nodes (\texttt{src/yuho/ast/nodes.py}). The frozen-dataclass
choice is deliberate: every analyser, transpiler, and diagnostic pass takes
the same AST as input, so structural sharing and equality-by-value matter.

\subsection{Six transpilers}
\label{subsec:transpilers}
Each transpiler is a Visitor over the AST.
\begin{description}
\item[JSON] structural dump for downstream tooling.
\item[Controlled English] template-driven natural-language rendering.
\item[\LaTeX{}] camera-ready statute prints with footnoted SSO anchors.
\item[Mermaid] decision-tree flowcharts for elements + exceptions.
\item[Alloy] a relational model for bounded model-checking
(\texttt{src/yuho/transpile/alloy\_transpiler.py}, $\sim$26k SLOC).
\item[DOCX] minimal OOXML \texttt{.docx} package built with stdlib
\texttt{zipfile} + WordprocessingML; no \texttt{python-docx} dependency.
\end{description}

\subsection{Verification: Z3 and Alloy}
\todo{1--2 paragraphs: what kinds of properties we can check today
(consistency, exception-priority well-formedness, mutual-exclusion of
penalty branches), what we cannot yet (deep precedent reasoning).}
Z3 powers exception-priority validation and conflict detection between
sections. Alloy emits relational models suitable for bounded model-checking;
the typical use case is exhaustive enumeration of element-combination
counterexamples.

\subsection{Three-tier coverage harness}
\label{subsec:coverage}
\begin{description}
\item[L1] tree-sitter parses without error.
\item[L2] AST builds and passes static semantic checks (lint, type, fidelity).
\item[L3] human review using a structured 11-point checklist
(illustration count vs canonical, penalty fidelity, element completeness,
amendment lineage, exception priority sanity, etc.).
\end{description}
The L3 reviewer is dispatcher-driven (\texttt{scripts/phase\_d\_l3\_review.py})
and supports parallel agent sessions stamping signoffs into per-section
\texttt{metadata.toml}.

\subsection{Editor surfaces}
\label{subsec:editors}
A pygls-based Language Server provides hover, inlay hints, code lenses,
diagnostics, and contextual completion. A Model Context Protocol server
exposes the same workflow as machine-callable tools, resources, and prompts
for AI clients (Claude, Codex CLI, Cursor). A VS Code extension
(\texttt{editors/vscode-yuho/}, v0.2) wraps the LSP and adds snippets, a
status-bar coverage indicator, and commands to run \texttt{yuho check} or
transpile in-editor.

\subsection{Fidelity diagnostics}
\label{subsec:fidelity}
\todo{Detail the four fidelity checks added in the diagnostics layer
(\textsf{G4} illustration-count, fabricated fine cap, fabricated caning
range, \textsf{G11} disjunctive-connective). Concrete example for each.}
The diagnostic pass cross-checks the encoded AST against
\texttt{library/penal\_code/\_raw/act.json} (the scraped canonical text)
to catch the most common fidelity errors: missing illustrations, fabricated
penalty caps, and conjunctive-vs-disjunctive miscoding.

\subsection{Engineering numbers}
\todo{Table: SLOC per layer (parser, AST, analysers, transpilers, CLI,
LSP, MCP, scraper). Auto-generated by \texttt{scripts/repo\_stats.py}.}
